\uuid{Gw2t}
\exo7id{3569}
\titre{exo7 3569}
\auteur{quercia}
\organisation{exo7}
\datecreate{2010-03-10}
\isIndication{false}
\isCorrection{false}
\chapitre{Endomorphisme particulier}
\sousChapitre{Autre}
\cpge{mpsi-mp2i.espaces-vectoriels.a, pcsi.espaces-vectoriels.a, ptsi.geometrie-espace.c, mpsi-mp2i.espaces-vectoriels.b, pcsi.espaces-vectoriels.b, ptsi.espaces-vectoriels.a, mpsi-mp2i.applications-lineaires.b, pcsi.applications-lineaires.b, ptsi.applications-lineaires.b, bcpst1.espace-vectoriel-kn.b, bcpst2.applications-lineaires-matrices.a}
\module{Algèbre}
\niveau{L2}
\difficulte{}

\contenu{
\question{
Soit $E$ un $ K$-espace vectoriel de dimension $n$, $(F,G)$ deux sous-espaces vectoriels
de $E$ tels que 
$E = F \oplus G$. On note $p$ la projection sur $F$ parallèlement à~$G$.
Soit $E_p = \{f\in\mathcal{L}(E) \text{ tq } f\circ p = p\circ f\}$.
Quelle est la dimension de $E_p$ ?
}
}
